\documentclass[aspectratio=169]{beamer}

% Theme académique professionnel
\usetheme{Madrid}
\usecolortheme{default}

% Packages
\usepackage{amsmath,amsfonts}
\usepackage[french]{babel}
\usepackage{tikz}
\usepackage{pgfplots}
\pgfplotsset{compat=1.18}



% Informations
\title{Time Series}
\subtitle{Stationnary Time Series}
\author{Mr. Aymen Ben Brik}

\institute{
Time Series Course 1GAMMA
}

\date{\today}


\usepackage{listings}
\usepackage{xcolor}

\lstset{
language=Python,
basicstyle=\ttfamily\small,
keywordstyle=\color{blue},
commentstyle=\color{gray},
stringstyle=\color{red},
showstringspaces=false,
frame=single,
breaklines=true
}



\begin{document}

% Slide titre
\begin{frame}
\titlepage
\end{frame}

% Objectifs
\begin{frame}{Session Objectives}
\begin{itemize}
    \item Understand the concept of stationarity in time series
    \item Distinguish between strict and weak stationarity
    \item Learn how to test for stationarity (ADF, KPSS)
    \item Interpret ACF and PACF plots
    \item Introduce AR, MA, and ARMA models
    \item Diagnose model adequacy through residual analysis
\end{itemize}
\end{frame}

\section{Les séries temporelles stationnaires}

\subsection{La stationnarité}
\begin{frame}{What is a Stationary Time Series?}
\begin{itemize}
    \item A time series is said to be \textbf{stationary} if its statistical properties do not change over time
    \vspace{0.3cm}
    \item In particular, a stationary series has:
    \begin{itemize}
        \item Constant mean
        \item Constant variance
        \item Constant autocovariance (depends only on the lag)
    \end{itemize}
    \vspace{0.3cm}
    \item Intuitively, the series fluctuates around a stable level with no systematic trend or changing variability
    \vspace{0.3cm}
    \item Stationarity is a key assumption for many time series models (AR, MA, ARMA)
\end{itemize}
\end{frame}

\begin{frame}{Stationary vs Non-Stationary Series}
\begin{itemize}
    \item \textbf{Stationary series:}
    \begin{itemize}
        \item Fluctuates around a constant mean
        \item Variance remains stable over time
        \item No clear trend or seasonality
    \end{itemize}
    
    \vspace{0.4cm}
    
    \item \textbf{Non-stationary series:}
    \begin{itemize}
        \item Presence of trend (upward or downward)
        \item Changing variance over time
        \item Possible seasonality or structural changes
    \end{itemize}
    
    \vspace{0.4cm}
    
    \item \textbf{Examples:}
    \begin{itemize}
        \item Stationary: white noise
        \item Non-stationary: random walk, trending data
    \end{itemize}
\end{itemize}
\end{frame}

\begin{frame}{Visual Illustration: Stationary vs Non-Stationary}
\centering
\begin{tikzpicture}
\begin{axis}[
    width=0.9\textwidth,
    height=0.65\textheight,
    xlabel={Time},
    ylabel={Value},
    legend style={at={(0.5,-0.2)}, anchor=north, legend columns=2},
    ymin=-1, ymax=8,
    xmin=0, xmax=10,
    samples=100,
    domain=0:10,
    grid=both
]

% Stationary series
\addplot[blue, thick] {2 + 0.4*sin(deg(4*x)) + 0.2*sin(deg(11*x))};
\addlegendentry{Stationary series}

% Non-stationary series
\addplot[red, thick] {0.5*x + 0.6*sin(deg(3*x)) + 1};
\addlegendentry{Non-stationary series}

\end{axis}
\end{tikzpicture}

\vspace{0.2cm}
\begin{itemize}
    \item The \textcolor{blue}{blue series} fluctuates around a constant mean
    \item The \textcolor{red}{red series} shows an upward trend over time
\end{itemize}
\end{frame}

\begin{frame}[fragile]{Simulation of a White Noise Process}
\small
\textbf{Definition:} A white noise process is a sequence of uncorrelated random variables with zero mean and constant variance.

\vspace{0.2cm}

\begin{columns}[T]
\begin{column}{0.48\textwidth}
\textbf{R}
\begin{verbatim}
set.seed(123)
n <- 200
wn <- rnorm(n, 0, 1)
plot.ts(wn, col="blue",
 main="White Noise in R")
\end{verbatim}
\end{column}

\begin{column}{0.48\textwidth}
\textbf{Python}
\begin{verbatim}
import numpy as np
import matplotlib.pyplot as plt
np.random.seed(123)
wn = np.random.normal(0, 1, 200)
plt.plot(wn)
plt.title("White Noise in Python")
plt.show()
\end{verbatim}
\end{column}
\end{columns}

\vspace{0.1cm}
\textbf{Observation:} The series fluctuates randomly around zero.
\end{frame}

\begin{frame}[fragile]{Random Walk Simulation}
\small
\textbf{Model:}
\[
X_t = X_{t-1} + \varepsilon_t, \qquad \varepsilon_t \sim WN(0,\sigma^2)
\]

\vspace{0.2cm}

\begin{columns}[T]
\begin{column}{0.48\textwidth}
\begin{block}{R}
\begin{verbatim}
set.seed(123)
n <- 200
eps <- rnorm(n, 0, 1)
rw <- cumsum(eps)

plot.ts(rw, col="red",
 main="Random Walk (R)")
\end{verbatim}
\end{block}
\end{column}

\begin{column}{0.48\textwidth}
\begin{block}{Python}
\begin{verbatim}
import numpy as np
import matplotlib.pyplot as plt

np.random.seed(123)
eps = np.random.normal(0, 1, 200)
rw = np.cumsum(eps)

plt.plot(rw)
plt.title("Random Walk (Python)")
plt.show()
\end{verbatim}
\end{block}
\end{column}
\end{columns}

\vspace{0.2cm}
\textbf{Key idea:} The accumulation of shocks makes the process non-stationary.
\end{frame}



\begin{frame}{Strong Stationarity}
\textbf{Definition:}

\vspace{0.3cm}

A time series $\{X_t\}$ is said to be \textbf{strictly (strongly) stationary} if:

\begin{itemize}
    \item For any integer $k$, any time points $t_1, \dots, t_k$, and any shift $h$,
    \item The joint distribution of 
    \[
    (X_{t_1}, X_{t_2}, \dots, X_{t_k})
    \]
    is identical to that of
    \[
    (X_{t_1+h}, X_{t_2+h}, \dots, X_{t_k+h})
    \]
\end{itemize}

\vspace{0.4cm}

\textbf{Key idea:}
\begin{itemize}
    \item The probabilistic behavior of the series does not change over time
\end{itemize}

\end{frame}

\begin{frame}{Intuition of Strong Stationarity}

\textbf{Interpretation:}
\begin{itemize}
    \item The entire probabilistic structure of the series is invariant over time
    \item Not only mean and variance, but \textbf{all distributions} remain unchanged
\end{itemize}

\vspace{0.4cm}

\textbf{Example (Stationary):}
\begin{itemize}
    \item White noise process:
    \[
    X_t \sim \mathcal{N}(0, \sigma^2)
    \]
    \item Same distribution at all times
\end{itemize}

\vspace{0.4cm}

\textbf{Counterexample (Non-Stationary):}
\begin{itemize}
    \item Time-varying variance:
    \[
    X_t \sim \mathcal{N}(0, t)
    \]
    \item Distribution changes over time
\end{itemize}

\vspace{0.4cm}

\textbf{Key takeaway:}
\begin{itemize}
    \item Strong stationarity is a very strict condition, rarely verified in practice
\end{itemize}

\end{frame}

\begin{frame}{Weak Stationarity (Covariance Stationarity)}

\textbf{Definition:}

\vspace{0.3cm}

A time series $\{X_t\}$ is said to be \textbf{weakly stationary} if:

\begin{itemize}
    \item $\mathbb{E}[X_t] = \mu$ (constant mean)
    \item $\text{Var}(X_t) = \sigma^2$ (constant variance)
    \item $\text{Cov}(X_t, X_{t+h}) = \gamma(h)$ depends only on the lag $h$
\end{itemize}

\vspace{0.4cm}

\textbf{Interpretation:}
\begin{itemize}
    \item The first two moments are time-invariant
    \item The dependence structure depends only on the time gap, not on time itself
\end{itemize}

\vspace{0.4cm}

\textbf{Key insight:}
\begin{itemize}
    \item Weak stationarity is sufficient for most models (AR, MA, ARMA)
\end{itemize}

\end{frame}

\begin{frame}{Strong vs Weak Stationarity}

\begin{center}
\begin{tabular}{|c|c|c|}
\hline
\textbf{Property} & \textbf{Strong Stationarity} & \textbf{Weak Stationarity} \\
\hline
Distribution & Unchanged over time & Not required \\
\hline
Mean & Constant & Constant \\
\hline
Variance & Constant & Constant \\
\hline
Autocovariance & Implicitly constant & Depends only on lag \\
\hline
Conditions & Very restrictive & Less restrictive \\
\hline
Practical use & Rare in practice & Widely used (AR, MA, ARMA) \\
\hline
\end{tabular}
\end{center}

\vspace{0.4cm}

\textbf{Key takeaway:}
\begin{itemize}
    \item Strong stationarity $\Rightarrow$ Weak stationarity (under finite variance)
    \item In practice, we usually assume \textbf{weak stationarity}
\end{itemize}

\end{frame}

\begin{frame}{How to Check Stationarity (Graphical Intuition)}

\textbf{Visual Inspection:}
\begin{itemize}
    \item Plot the time series and observe its behavior over time
\end{itemize}

\vspace{0.4cm}

\textbf{Signs of Stationarity:}
\begin{itemize}
    \item Fluctuates around a constant mean
    \item Constant variance (no increasing/decreasing spread)
    \item No visible trend or seasonality
\end{itemize}

\vspace{0.4cm}

\textbf{Signs of Non-Stationarity:}
\begin{itemize}
    \item Upward or downward trend
    \item Changing variance over time
    \item Presence of seasonal patterns
\end{itemize}

\vspace{0.4cm}

\textbf{Tools:}
\begin{itemize}
    \item Time series plot
    \item Rolling mean and rolling variance
\end{itemize}

\end{frame}

\begin{frame}[fragile]{Rolling Statistics (Mean and Variance)}

\small
\textbf{Idea: Compute statistics over a moving window}

\begin{columns}[T]
\begin{column}{0.48\textwidth}
\begin{block}{R}
\begin{verbatim}
library(zoo)

n <- 200
x <- cumsum(rnorm(n))  # example

roll_mean <- rollmean(x, k=20,
 fill=NA)
roll_var <- rollapply(x, 20,
 var, fill=NA)

plot(x, type="l")
lines(roll_mean, col="blue")
\end{verbatim}
\end{block}
\end{column}

\begin{column}{0.48\textwidth}
\begin{block}{Python}
\begin{verbatim}
import numpy as np
import pandas as pd
import matplotlib.pyplot as plt

x = np.cumsum(
 np.random.normal(0,1,200))

s = pd.Series(x)
roll_mean = s.rolling(20).mean()
roll_var = s.rolling(20).var()

plt.plot(s)
plt.plot(roll_mean)
plt.show()
\end{verbatim}
\end{block}
\end{column}
\end{columns}

\vspace{0.2cm}
\textbf{Interpretation:}
\begin{itemize}
    \item Constant rolling mean/variance → likely stationary
    \item Changing statistics → non-stationary
\end{itemize}

\end{frame}
\begin{frame}{Theoretical Exercise}

\textbf{Exercise:}

Consider the following time series:

\begin{enumerate}
    \item $X_t = \varepsilon_t$, where $\varepsilon_t \sim WN(0,\sigma^2)$
    \item $X_t = t + \varepsilon_t$
    \item $X_t = X_{t-1} + \varepsilon_t$
    \item $X_t = \varepsilon_t + 0.5\,\varepsilon_{t-1}$
\end{enumerate}

\vspace{0.4cm}

\textbf{Questions:}
\begin{itemize}
    \item Which of these processes are stationary?
    \item Justify your answer (mean, variance, dependence structure)
\end{itemize}

\vspace{0.4cm}
\end{frame}



\subsection{Les tests de stationnarité}

\begin{frame}{ADF Test (Augmented Dickey-Fuller)}

\textbf{Objective:}
\begin{itemize}
    \item Test whether a time series is \textbf{stationary}
\end{itemize}

\vspace{0.4cm}

\textbf{Key idea:}
\begin{itemize}
    \item Detect the presence of a \textbf{unit root}
\end{itemize}

\vspace{0.4cm}

\textbf{Why important?}
\begin{itemize}
    \item A unit root implies \textbf{non-stationarity}
    \item Many models (ARMA) require stationarity
\end{itemize}
\end{frame}

\begin{frame}{ADF Test: Model}

\textbf{Basic idea:}

\vspace{0.3cm}

Start from an AR(1) model:
\[
X_t = \rho X_{t-1} + \varepsilon_t
\]

\vspace{0.3cm}

Rewriting:
\[
\Delta X_t = \gamma X_{t-1} + \varepsilon_t
\quad \text{with } \gamma = \rho - 1
\]

\vspace{0.4cm}

\textbf{Augmented version:}
\[
\Delta X_t = \gamma X_{t-1} + \sum_{i=1}^{p} \alpha_i \Delta X_{t-i} + \varepsilon_t
\]

\vspace{0.3cm}

\textbf{Goal:} Test whether $\gamma = 0$

\end{frame}

\begin{frame}{ADF Test: Hypotheses}

\textbf{Null hypothesis:}
\begin{itemize}
    \item $H_0: \gamma = 0$
    \item The series has a \textbf{unit root}
    \item $\Rightarrow$ Non-stationary
\end{itemize}

\vspace{0.4cm}

\textbf{Alternative hypothesis:}
\begin{itemize}
    \item $H_1: \gamma < 0$
    \item The series is \textbf{stationary}
\end{itemize}

\vspace{0.4cm}

\textbf{Interpretation:}
\begin{itemize}
    \item Reject $H_0$ $\Rightarrow$ Stationary series
    \item Fail to reject $H_0$ $\Rightarrow$ Non-stationary
\end{itemize}

\end{frame}

\begin{frame}{ADF Test: Decision Rule}

\textbf{Based on p-value:}

\begin{itemize}
    \item p-value $< 0.05$ $\Rightarrow$ Reject $H_0$ $\Rightarrow$ Stationary
    \item p-value $> 0.05$ $\Rightarrow$ Fail to reject $H_0$ $\Rightarrow$ Non-stationary
\end{itemize}

\vspace{0.4cm}

\textbf{Important:}
\begin{itemize}
    \item The test is sensitive to:
    \begin{itemize}
        \item Lag selection
        \item Trend or intercept specification
    \end{itemize}
\end{itemize}

\end{frame}

\begin{frame}[fragile]{ADF Test: Implementation in R and Python}
\small

\begin{columns}[T]
\begin{column}{0.48\textwidth}
\begin{block}{R}
\begin{verbatim}
library(tseries)

x <- cumsum(rnorm(200))   # example series

adf.test(x)
\end{verbatim}

\vspace{0.2cm}
\textbf{Output:}
\begin{itemize}
    \item test statistic
    \item p-value
\end{itemize}
\end{block}
\end{column}

\begin{column}{0.48\textwidth}
\begin{block}{Python}
\begin{verbatim}
import numpy as np
from statsmodels.tsa.stattools import adfuller

x = np.cumsum(np.random.normal(0, 1, 200))

result = adfuller(x)

print("ADF Statistic:", result[0])
print("p-value:", result[1])
\end{verbatim}

\vspace{0.2cm}
\textbf{Output:}
\begin{itemize}
    \item ADF statistic
    \item p-value
\end{itemize}
\end{block}
\end{column}
\end{columns}

\end{frame}

\begin{frame}{ADF Test: Interpretation of Results}

\textbf{Decision rule based on the p-value:}
\begin{itemize}
    \item If p-value $< 0.05$: reject $H_0$
    \item If p-value $> 0.05$: fail to reject $H_0$
\end{itemize}

\vspace{0.4cm}

\textbf{Interpretation:}
\begin{itemize}
    \item Rejecting $H_0$ means that the series is likely \textbf{stationary}
    \item Failing to reject $H_0$ suggests the presence of a \textbf{unit root}
    \item Therefore, the series is likely \textbf{non-stationary}
\end{itemize}

\vspace{0.4cm}

\textbf{Reminder:}
\begin{itemize}
    \item $H_0$: the series has a unit root
    \item $H_1$: the series is stationary
\end{itemize}

\end{frame}


\begin{frame}{KPSS Test (Kwiatkowski-Phillips-Schmidt-Shin)}

\textbf{Objective:}
\begin{itemize}
    \item Test whether a time series is \textbf{stationary}
\end{itemize}

\vspace{0.4cm}

\textbf{Key difference from ADF:}
\begin{itemize}
    \item The KPSS test takes \textbf{stationarity} as the null hypothesis
    \item It is therefore complementary to the ADF test
\end{itemize}

\vspace{0.4cm}

\textbf{Why useful?}
\begin{itemize}
    \item Using both ADF and KPSS provides a more reliable conclusion
\end{itemize}

\end{frame}

\begin{frame}{KPSS Test: Hypotheses}

\textbf{Null hypothesis:}
\begin{itemize}
    \item $H_0$: the series is \textbf{stationary}
\end{itemize}

\vspace{0.3cm}

\textbf{Alternative hypothesis:}
\begin{itemize}
    \item $H_1$: the series is \textbf{non-stationary}
\end{itemize}

\vspace{0.4cm}

\textbf{Interpretation:}
\begin{itemize}
    \item Reject $H_0$ $\Rightarrow$ evidence of non-stationarity
    \item Fail to reject $H_0$ $\Rightarrow$ the series may be stationary
\end{itemize}

\end{frame}

\begin{frame}{ADF vs KPSS}

\begin{center}
\begin{tabular}{|c|c|c|}
\hline
\textbf{Test} & \textbf{Null Hypothesis} & \textbf{Alternative} \\
\hline
ADF & Unit root / Non-stationary & Stationary \\
\hline
KPSS & Stationary & Non-stationary \\
\hline
\end{tabular}
\end{center}

\vspace{0.4cm}

\textbf{Practical interpretation:}
\begin{itemize}
    \item ADF rejects + KPSS does not reject $\Rightarrow$ likely stationary
    \item ADF does not reject + KPSS rejects $\Rightarrow$ likely non-stationary
\end{itemize}

\end{frame}

\begin{frame}[fragile]{KPSS Test: Implementation in R and Python}
\small

\begin{columns}[T]
\begin{column}{0.48\textwidth}
\begin{block}{R}
\begin{verbatim}
library(tseries)

x <- cumsum(rnorm(200))   # example series

kpss.test(x)
\end{verbatim}
\end{block}
\end{column}

\begin{column}{0.48\textwidth}
\begin{block}{Python}
\begin{verbatim}
import numpy as np
from statsmodels.tsa.stattools import kpss

x = np.cumsum(np.random.normal(0, 1, 200))

stat, p_value, _, _ = kpss(x)

print("KPSS Statistic:", stat)
print("p-value:", p_value)
\end{verbatim}
\end{block}
\end{column}
\end{columns}

\end{frame}

\subsection{ACF / PACF}

\begin{frame}{ACF and PACF}

\textbf{Objective:}
\begin{itemize}
    \item Analyze the dependence structure of a time series
\end{itemize}

\vspace{0.4cm}

\textbf{Key tools:}
\begin{itemize}
    \item ACF: Autocorrelation Function
    \item PACF: Partial Autocorrelation Function
\end{itemize}

\vspace{0.4cm}

\textbf{Why important?}
\begin{itemize}
    \item Identify the structure of time series models
    \item Help determine AR and MA orders
\end{itemize}

\end{frame}

\begin{frame}{Autocorrelation Function (ACF)}

\textbf{Definition:}
\begin{itemize}
    \item The ACF measures the correlation between $X_t$ and $X_{t+h}$
\end{itemize}

\vspace{0.3cm}

\[
\rho(h) = \frac{\mathrm{Cov}(X_t, X_{t+h})}{\mathrm{Var}(X_t)}
\]

\vspace{0.4cm}

\textbf{Interpretation:}
\begin{itemize}
    \item Measures linear dependence at lag $h$
    \item Values between $-1$ and $1$
\end{itemize}

\vspace{0.4cm}

\textbf{Key idea:}
\begin{itemize}
    \item High autocorrelation $\Rightarrow$ strong dependence
\end{itemize}

\end{frame}

\begin{frame}{ACF: Intuition}

\textbf{What does ACF tell us?}

\begin{itemize}
    \item How past values influence current values
\end{itemize}

\vspace{0.4cm}

\textbf{Examples:}
\begin{itemize}
    \item White noise:
    \begin{itemize}
        \item No correlation $\Rightarrow$ ACF $\approx 0$ for all lags
    \end{itemize}
    
    \vspace{0.2cm}
    
    \item Persistent series:
    \begin{itemize}
        \item Slowly decreasing ACF
    \end{itemize}
\end{itemize}

\vspace{0.4cm}

\textbf{Graph:}
\begin{itemize}
    \item ACF is typically represented as a bar plot
\end{itemize}

\end{frame}

\begin{frame}{Partial Autocorrelation Function (PACF)}

\textbf{Definition:}
\begin{itemize}
    \item The PACF measures the direct correlation between $X_t$ and $X_{t+h}$
    \item Removing the effect of intermediate lags
\end{itemize}

\vspace{0.4cm}

\textbf{Key idea:}
\begin{itemize}
    \item PACF isolates the \textbf{pure effect} of lag $h$
\end{itemize}

\vspace{0.4cm}

\textbf{Interpretation:}
\begin{itemize}
    \item Helps identify the order of AR models
\end{itemize}

\end{frame}

\begin{frame}{ACF vs PACF}

\begin{center}
\begin{tabular}{|c|c|c|}
\hline
\textbf{Function} & \textbf{Measures} & \textbf{Use} \\
\hline
ACF & Total correlation & Identify MA(q) \\
\hline
PACF & Direct correlation & Identify AR(p) \\
\hline
\end{tabular}
\end{center}

\vspace{0.4cm}

\textbf{Key idea:}
\begin{itemize}
    \item ACF includes indirect effects
    \item PACF removes indirect effects
\end{itemize}

\end{frame}

\begin{frame}[fragile]{Simulation of an AR(2): ACF and PACF}
\small
\textbf{AR(2) model:}
\[
X_t = 0.6X_{t-1} - 0.3X_{t-2} + \varepsilon_t,
\qquad \varepsilon_t \sim WN(0,1)
\]

\begin{columns}[T]
\begin{column}{0.48\textwidth}
\begin{block}{R}
\begin{verbatim}
set.seed(123)
x <- arima.sim(n = 200,
 model = list(ar = c(0.6, -0.3)))

par(mfrow = c(1,2))
acf(x, main = "ACF of AR(2)")
pacf(x, main = "PACF of AR(2)")
\end{verbatim}
\end{block}
\end{column}

\begin{column}{0.48\textwidth}
\begin{block}{Python}
\begin{verbatim}
import numpy as np
from statsmodels.tsa.arima_process import ArmaProcess
from statsmodels.graphics.tsaplots import plot_acf, plot_pacf
import matplotlib.pyplot as plt

np.random.seed(123)
ar = np.array([1, -0.6, 0.3])
ma = np.array([1])
x = ArmaProcess(ar, ma).generate_sample(200)

fig, ax = plt.subplots(1, 2, figsize=(10,4))
plot_acf(x, ax=ax[0], lags=20)
plot_pacf(x, ax=ax[1], lags=20, method="ywm")
plt.show()
\end{verbatim}
\end{block}
\end{column}
\end{columns}
\end{frame}

\begin{frame}{Interpretation of ACF and PACF for an AR(2)}

\textbf{Typical pattern of an AR(2):}
\begin{itemize}
    \item The \textbf{ACF} decreases gradually
    \item The \textbf{PACF} cuts off after lag 2
\end{itemize}

\vspace{0.4cm}

\textbf{Why?}
\begin{itemize}
    \item In an AR(p) process, the PACF truncates at lag $p$
    \item For an AR(2), only lags 1 and 2 have direct effects
\end{itemize}

\vspace{0.4cm}

\textbf{Key takeaway:}
\begin{itemize}
    \item Slowly decaying ACF + PACF cutoff at lag 2
    \item This is the signature of an \textbf{AR(2)} process
\end{itemize}

\end{frame}

\begin{frame}[fragile]{Simulation of an MA(3): ACF and PACF}
\small
\textbf{MA(3) model:}
\[
X_t = \varepsilon_t + 0.5\,\varepsilon_{t-1} + 0.3\,\varepsilon_{t-2} + 0.2\,\varepsilon_{t-3},
\qquad \varepsilon_t \sim WN(0,1)
\]
\begin{columns}[T]
\begin{column}{0.48\textwidth}
\begin{block}{R}
\begin{verbatim}
set.seed(123)
x <- arima.sim(n = 200,
 model = list(ma = c(0.5, 0.3, 0.2)))

par(mfrow = c(1,2))
acf(x, main = "ACF of MA(3)")
pacf(x, main = "PACF of MA(3)")
\end{verbatim}
\end{block}
\end{column}

\begin{column}{0.48\textwidth}
\begin{block}{Python}
\begin{verbatim}
import numpy as np
from statsmodels.tsa.arima_process import ArmaProcess
from statsmodels.graphics.tsaplots import plot_acf, plot_pacf
import matplotlib.pyplot as plt

np.random.seed(123)
ar = np.array([1])
ma = np.array([1, 0.5, 0.3, 0.2])
x = ArmaProcess(ar, ma).generate_sample(200)

fig, ax = plt.subplots(1, 2, figsize=(10,4))
plot_acf(x, ax=ax[0], lags=20)
plot_pacf(x, ax=ax[1], lags=20, method="ywm")
plt.show()
\end{verbatim}
\end{block}
\end{column}
\end{columns}

\end{frame}

\begin{frame}{Interpretation of ACF and PACF for an MA(3)}

\textbf{Typical pattern of an MA(3):}
\begin{itemize}
    \item The \textbf{ACF} cuts off after lag 3
    \item The \textbf{PACF} decreases gradually
\end{itemize}

\vspace{0.4cm}

\textbf{Why?}
\begin{itemize}
    \item In an MA(q) process, the ACF truncates at lag $q$
    \item Dependence exists only through past shocks up to lag 3
\end{itemize}

\vspace{0.4cm}

\textbf{Key takeaway:}
\begin{itemize}
    \item ACF cutoff + PACF decay $\Rightarrow$ MA process
    \item Here: cutoff at lag 3 $\Rightarrow$ MA(3)
\end{itemize}

\end{frame}

\begin{frame}{Summary: ACF/PACF Patterns (AR vs MA vs ARMA)}

\begin{center}
\begin{tabular}{|c|c|c|}
\hline
\textbf{Model} & \textbf{ACF} & \textbf{PACF} \\
\hline
AR(p) & Gradual decay & Cuts off after lag $p$ \\
\hline
MA(q) & Cuts off after lag $q$ & Gradual decay \\
\hline
ARMA(p,q) & Gradual decay & Gradual decay \\
\hline
\end{tabular}
\end{center}

\vspace{0.4cm}

\textbf{Key identification rules:}
\begin{itemize}
    \item Sharp cutoff in PACF $\Rightarrow$ AR model
    \item Sharp cutoff in ACF $\Rightarrow$ MA model
    \item No clear cutoff $\Rightarrow$ ARMA model
\end{itemize}

\vspace{0.4cm}

\textbf{Important:}
\begin{itemize}
    \item In practice, patterns are not always perfectly clear
    \item Use ACF/PACF together with model selection criteria (AIC, BIC)
\end{itemize}

\end{frame}

\subsection{Les modèles AR / MA / ARMA}
\begin{frame}{Autoregressive Model AR(p)}

\textbf{Definition:}

\vspace{0.3cm}

An autoregressive process of order $p$, denoted AR($p$), is defined as:

\[
X_t = \phi_1 X_{t-1} + \phi_2 X_{t-2} + \cdots + \phi_p X_{t-p} + \varepsilon_t
\]

\vspace{0.3cm}

where:
\begin{itemize}
    \item $\phi_1, \dots, \phi_p$ are parameters
    \item $\varepsilon_t \sim WN(0,\sigma^2)$
\end{itemize}

\vspace{0.4cm}

\textbf{Key idea:}
\begin{itemize}
    \item The current value depends on its \textbf{past values}
\end{itemize}

\end{frame}

\begin{frame}{AR(p): Intuition}

\textbf{How does it work?}

\begin{itemize}
    \item The present is a linear combination of the past
\end{itemize}

\vspace{0.4cm}

\textbf{Example: AR(1)}
\[
X_t = \phi X_{t-1} + \varepsilon_t
\]

\begin{itemize}
    \item If $|\phi| < 1$: the process is stationary
    \item The series tends to return to its mean (mean reversion)
\end{itemize}

\vspace{0.4cm}

\textbf{Interpretation:}
\begin{itemize}
    \item Large $\phi$ $\Rightarrow$ strong persistence
    \item Small $\phi$ $\Rightarrow$ weak dependence
\end{itemize}

\end{frame}

\begin{frame}{AR(p): Example}

\textbf{Example: AR(2)}

\[
X_t = 0.6 X_{t-1} - 0.3 X_{t-2} + \varepsilon_t
\]

\vspace{0.4cm}

\textbf{Interpretation:}
\begin{itemize}
    \item $X_{t-1}$ has a positive influence
    \item $X_{t-2}$ has a negative influence
\end{itemize}

\vspace{0.4cm}

\textbf{Key idea:}
\begin{itemize}
    \item The structure captures dynamic dependence over time
\end{itemize}

\end{frame}

\begin{frame}{Stationarity Condition of AR(p)}

\textbf{AR($p$) model:}
\[
X_t = \phi_1 X_{t-1} + \cdots + \phi_p X_{t-p} + \varepsilon_t
\]

\vspace{0.4cm}

\textbf{Characteristic polynomial:}
\[
\Phi(z) = 1 - \phi_1 z - \phi_2 z^2 - \cdots - \phi_p z^p
\]

\vspace{0.4cm}

\textbf{Stationarity condition:}
\begin{itemize}
    \item The process is stationary if all roots of $\Phi(z)=0$ lie \textbf{outside the unit circle}
\end{itemize}

\vspace{0.4cm}

\textbf{Equivalent condition:}
\begin{itemize}
    \item All roots satisfy $|z_i| > 1$
\end{itemize}

\end{frame}

\begin{frame}{Stationarity Condition: AR(1)}

\textbf{Model:}
\[
X_t = \phi X_{t-1} + \varepsilon_t
\]

\vspace{0.4cm}

\textbf{Characteristic equation:}
\[
1 - \phi z = 0 \quad \Rightarrow \quad z = \frac{1}{\phi}
\]

\vspace{0.4cm}

\textbf{Stationarity condition:}
\[
|z| > 1 \quad \Rightarrow \quad \left|\frac{1}{\phi}\right| > 1
\]

\[
\Rightarrow \quad |\phi| < 1
\]

\vspace{0.4cm}

\textbf{Interpretation:}
\begin{itemize}
    \item If $|\phi| < 1$: stationary (mean-reverting)
    \item If $|\phi| \geq 1$: non-stationary
\end{itemize}

\end{frame}

\begin{frame}{Moving Average Model MA(q)}

\textbf{Definition:}

\vspace{0.3cm}

A moving average process of order $q$, denoted MA($q$), is defined as:

\[
X_t = \varepsilon_t + \theta_1 \varepsilon_{t-1} + \theta_2 \varepsilon_{t-2} + \cdots + \theta_q \varepsilon_{t-q}
\]

\vspace{0.3cm}

where:
\begin{itemize}
    \item $\varepsilon_t \sim WN(0,\sigma^2)$
    \item $\theta_1,\dots,\theta_q$ are constants
\end{itemize}

\vspace{0.4cm}

\textbf{Key idea:}
\begin{itemize}
    \item The current value depends on the \textbf{current and past shocks}
\end{itemize}

\end{frame}

\begin{frame}{MA(q): Intuition}

\textbf{How does it work?}

\begin{itemize}
    \item A shock affects the series immediately
    \item Its effect persists only for a limited number of periods
\end{itemize}

\vspace{0.4cm}

\textbf{Interpretation:}
\begin{itemize}
    \item In an MA($q$) process, a shock at time $t$ influences:
    \begin{itemize}
        \item $X_t$
        \item $X_{t+1}$
        \item \dots
        \item $X_{t+q}$
    \end{itemize}
    \item After lag $q$, the effect disappears
\end{itemize}

\vspace{0.4cm}

\textbf{Key difference from AR(p):}
\begin{itemize}
    \item AR: depends on past values
    \item MA: depends on past shocks
\end{itemize}

\end{frame}

\begin{frame}{MA(q): Example}

\textbf{Example: MA(2)}

\[
X_t = \varepsilon_t + 0.6\,\varepsilon_{t-1} + 0.3\,\varepsilon_{t-2}
\]

\vspace{0.4cm}

\textbf{Interpretation:}
\begin{itemize}
    \item The current shock $\varepsilon_t$ affects $X_t$
    \item The previous shock $\varepsilon_{t-1}$ still affects $X_t$
    \item The shock $\varepsilon_{t-2}$ also affects $X_t$
    \item Older shocks have no effect
\end{itemize}

\vspace{0.4cm}

\textbf{Key idea:}
\begin{itemize}
    \item The memory of the process is finite
\end{itemize}

\end{frame}

\begin{frame}{Stationarity of MA(q)}

\textbf{Important property:}
\begin{itemize}
    \item Every MA($q$) process is \textbf{stationary}
\end{itemize}

\vspace{0.4cm}

\textbf{Why?}
\begin{itemize}
    \item $X_t$ is a finite linear combination of white noise terms
    \item Its mean is constant
    \item Its variance is constant
    \item Its autocovariance depends only on the lag
\end{itemize}

\vspace{0.4cm}

\textbf{Conclusion:}
\begin{itemize}
    \item MA models are always weakly stationary
\end{itemize}

\end{frame}
\begin{frame}{Mean and Variance of an MA(q)}

For
\[
X_t = \varepsilon_t + \theta_1\varepsilon_{t-1} + \cdots + \theta_q\varepsilon_{t-q},
\qquad \varepsilon_t \sim WN(0,\sigma^2),
\]

\vspace{0.3cm}

\textbf{Mean:}
\[
\mathbb{E}[X_t]=0
\]

\vspace{0.3cm}

\textbf{Variance:}
\[
\mathrm{Var}(X_t)=\sigma^2\left(1+\theta_1^2+\theta_2^2+\cdots+\theta_q^2\right)
\]

\vspace{0.4cm}

\textbf{Remark:}
\begin{itemize}
    \item The variance is constant over time
\end{itemize}

\end{frame}
\begin{frame}{Autocovariance of MA(q)}

\textbf{Model:}
\[
X_t = \varepsilon_t + \theta_1 \varepsilon_{t-1} + \cdots + \theta_q \varepsilon_{t-q},
\qquad \varepsilon_t \sim WN(0,\sigma^2)
\]

\vspace{0.4cm}

\textbf{Autocovariance function $\gamma(h)$:}

\begin{itemize}
    \item For $h = 0$:
    \[
    \gamma(0) = \sigma^2 \left(1 + \theta_1^2 + \cdots + \theta_q^2 \right)
    \]
    
    \vspace{0.3cm}
    
    \item For $1 \leq h \leq q$:
    \[
    \gamma(h) = \sigma^2 \left( \theta_h + \theta_1 \theta_{h+1} + \cdots + \theta_{q-h} \theta_q \right)
    \]
    
    \vspace{0.3cm}
    
    \item For $h > q$:
    \[
    \gamma(h) = 0
    \]
\end{itemize}

\vspace{0.4cm}

\textbf{Key property:}
\begin{itemize}
    \item The autocovariance function cuts off after lag $q$
\end{itemize}

\end{frame}

\begin{frame}{Invertibility of MA(q)}

\textbf{Model:}
\[
X_t = \varepsilon_t + \theta_1 \varepsilon_{t-1} + \cdots + \theta_q \varepsilon_{t-q}
\]

\vspace{0.4cm}

\textbf{Objective:}
\begin{itemize}
    \item Express $\varepsilon_t$ as a function of past values of $X_t$
\end{itemize}

\vspace{0.4cm}

\textbf{Key idea:}
\begin{itemize}
    \item An MA(q) process is \textbf{invertible} if it can be written as an AR($\infty$) process
\end{itemize}

\end{frame}


\begin{frame}{ARMA(p,q) Model}

\textbf{Definition:}

\vspace{0.3cm}

An ARMA($p,q$) process combines autoregressive and moving average components:

\[
X_t = \phi_1 X_{t-1} + \cdots + \phi_p X_{t-p}
+ \varepsilon_t + \theta_1 \varepsilon_{t-1} + \cdots + \theta_q \varepsilon_{t-q}
\]

\vspace{0.3cm}

where:
\begin{itemize}
    \item $\varepsilon_t \sim WN(0,\sigma^2)$
    \item $\phi_i$: AR parameters
    \item $\theta_j$: MA parameters
\end{itemize}

\end{frame}

\begin{frame}{ARMA(p,q): Compact Representation}

\textbf{Using the lag operator $B$:}

\vspace{0.3cm}

\[
\Phi(B)X_t = \Theta(B)\varepsilon_t
\]

\vspace{0.3cm}

where:

\[
\Phi(B) = 1 - \phi_1 B - \cdots - \phi_p B^p
\]

\[
\Theta(B) = 1 + \theta_1 B + \cdots + \theta_q B^q
\]

\vspace{0.4cm}

\textbf{Key idea:}
\begin{itemize}
    \item ARMA combines dependence on past values and past shocks
\end{itemize}

\end{frame}

\begin{frame}{ARMA(p,q): Conditions}

\textbf{Stationarity:}
\begin{itemize}
    \item All roots of $\Phi(z)=0$ lie outside the unit circle
\end{itemize}

\vspace{0.4cm}

\textbf{Invertibility:}
\begin{itemize}
    \item All roots of $\Theta(z)=0$ lie outside the unit circle
\end{itemize}

\vspace{0.4cm}

\textbf{Conclusion:}
\begin{itemize}
    \item ARMA is well-defined when both conditions are satisfied
\end{itemize}

\end{frame}

\begin{frame}{ARMA(p,q): Intuition}

\textbf{How does it work?}

\begin{itemize}
    \item AR part: captures persistence (dependence on past values)
    \item MA part: captures shocks (dependence on past errors)
\end{itemize}

\vspace{0.4cm}

\textbf{Key idea:}
\begin{itemize}
    \item ARMA provides a flexible model for time series dynamics
\end{itemize}

\vspace{0.4cm}

\textbf{Special cases:}
\begin{itemize}
    \item ARMA($p,0$) = AR($p$)
    \item ARMA($0,q$) = MA($q$)
\end{itemize}

\end{frame}


\begin{frame}{ARMA(p,q): ACF and PACF}

\textbf{Typical behavior:}

\begin{itemize}
    \item ACF: gradual decay
    \item PACF: gradual decay
\end{itemize}

\vspace{0.4cm}

\textbf{Key implication:}
\begin{itemize}
    \item No clear cutoff in ACF or PACF
    \item Harder to identify than AR or MA models
\end{itemize}

\vspace{0.4cm}

\textbf{In practice:}
\begin{itemize}
    \item Use AIC / BIC for model selection
\end{itemize}

\end{frame}


\subsection{Analyse des résidus}

\begin{frame}{Residual Analysis}

\textbf{Objective:}
\begin{itemize}
    \item Check whether the fitted model is adequate
\end{itemize}

\vspace{0.4cm}

\textbf{Key idea:}
\begin{itemize}
    \item The residuals should behave like \textbf{white noise}
\end{itemize}

\vspace{0.4cm}

\textbf{Why?}
\begin{itemize}
    \item If residuals contain structure, the model is not sufficient
\end{itemize}

\end{frame}

\begin{frame}{Properties of Good Residuals}

\textbf{A good model should produce residuals that:}

\begin{itemize}
    \item Have zero mean
    \item Have constant variance
    \item Are uncorrelated (no autocorrelation)
\end{itemize}

\vspace{0.4cm}

\textbf{In other words:}
\begin{itemize}
    \item Residuals $\sim$ White Noise
\end{itemize}

\end{frame}

\begin{frame}{Residual Diagnostics: Graphical Tools}

\textbf{Common plots:}
\begin{itemize}
    \item Residual time series plot
    \item Histogram of residuals
    \item Q-Q plot (normality check)
    \item ACF of residuals
\end{itemize}

\vspace{0.4cm}

\textbf{What to check:}
\begin{itemize}
    \item No visible pattern
    \item No trend or seasonality
    \item No significant autocorrelation
\end{itemize}

\end{frame}

\begin{frame}{ACF of Residuals}

\textbf{Key idea:}
\begin{itemize}
    \item Residuals should have no autocorrelation
\end{itemize}

\vspace{0.4cm}

\textbf{Interpretation:}
\begin{itemize}
    \item All autocorrelations should lie within confidence bounds
    \item Significant spikes $\Rightarrow$ model misspecification
\end{itemize}

\vspace{0.4cm}

\textbf{Conclusion:}
\begin{itemize}
    \item If ACF shows structure, the model is not adequate
\end{itemize}

\end{frame}


\begin{frame}{Ljung-Box Test}

\textbf{Objective:}
\begin{itemize}
    \item Test whether residuals are uncorrelated
\end{itemize}

\vspace{0.4cm}

\textbf{Hypotheses:}
\begin{itemize}
    \item $H_0$: residuals are independent (white noise)
    \item $H_1$: residuals are autocorrelated
\end{itemize}

\vspace{0.4cm}

\textbf{Decision:}
\begin{itemize}
    \item p-value $< 0.05$ $\Rightarrow$ reject $H_0$ (bad model)
    \item p-value $> 0.05$ $\Rightarrow$ good model
\end{itemize}

\end{frame}

\begin{frame}[fragile]{Residual Diagnostics in R}
\small

\begin{block}{R}
\begin{verbatim}
model <- arima(x, order=c(1,0,1))

res <- residuals(model)

acf(res)
Box.test(res, type="Ljung")
\end{verbatim}
\end{block}

\vspace{0.3cm}

\textbf{Goal:}
\begin{itemize}
    \item Check if residuals behave like white noise
\end{itemize}

\end{frame}

\begin{frame}[fragile]{Residual Diagnostics in Python}
\small

\begin{block}{Python}
\begin{verbatim}
from statsmodels.tsa.arima.model import ARIMA
from statsmodels.stats.diagnostic import acorr_ljungbox

model = ARIMA(x, order=(1,0,1)).fit()
res = model.resid

acorr_ljungbox(res, lags=[10])
\end{verbatim}
\end{block}

\vspace{0.3cm}

\textbf{Goal:}
\begin{itemize}
    \item Test residual independence using Ljung-Box
\end{itemize}

\end{frame}



\begin{frame}{Model Validation Checklist}

\textbf{A good model satisfies:}

\begin{itemize}
    \item Residuals $\sim$ white noise
    \item No autocorrelation (ACF flat)
    \item Ljung-Box test not rejected
\end{itemize}

\vspace{0.4cm}

\textbf{If not:}
\begin{itemize}
    \item Try a different model (change $p$, $q$)
    \item Consider transformations or differencing
\end{itemize}

\end{frame}




\subsection{Forecasting}

\begin{frame}{Forecasting in Time Series}

\textbf{Objective:}
\begin{itemize}
    \item Predict future values of a time series
\end{itemize}

\vspace{0.4cm}

\textbf{Key idea:}
\begin{itemize}
    \item Use past observations and a fitted model (AR, MA, ARMA)
\end{itemize}

\vspace{0.4cm}

\textbf{Applications:}
\begin{itemize}
    \item Economics (GDP, inflation)
    \item Finance (stock prices)
    \item Energy demand
\end{itemize}

\end{frame}

\begin{frame}{One-Step Ahead Forecast}

\textbf{Definition:}

\[
\hat{X}_{t+1} = \mathbb{E}[X_{t+1} \mid X_t, X_{t-1}, \dots]
\]

\vspace{0.4cm}

\textbf{Example: AR(1)}
\[
X_t = \phi X_{t-1} + \varepsilon_t
\]

\[
\Rightarrow \hat{X}_{t+1} = \phi X_t
\]

\vspace{0.4cm}

\textbf{Key idea:}
\begin{itemize}
    \item Future noise has zero expectation
\end{itemize}

\end{frame}

\begin{frame}{Multi-Step Forecast}

\textbf{Example: AR(1)}

\[
\hat{X}_{t+2} = \phi \hat{X}_{t+1} = \phi^2 X_t
\]

\[
\hat{X}_{t+h} = \phi^h X_t
\]

\vspace{0.4cm}

\textbf{Key idea:}
\begin{itemize}
    \item Forecasts converge to the mean when $|\phi|<1$
\end{itemize}

\vspace{0.4cm}

\textbf{Interpretation:}
\begin{itemize}
    \item The further ahead, the less precise the forecast
\end{itemize}

\end{frame}

\begin{frame}{Forecast Uncertainty}

\textbf{Important:}
\begin{itemize}
    \item Forecasts are not exact values
\end{itemize}

\vspace{0.4cm}

\textbf{Prediction intervals:}
\begin{itemize}
    \item Provide a range of possible future values
\end{itemize}

\vspace{0.4cm}

\textbf{Key idea:}
\begin{itemize}
    \item Uncertainty increases with the forecast horizon
\end{itemize}

\vspace{0.4cm}

\textbf{Example:}
\begin{itemize}
    \item Confidence intervals become wider over time
\end{itemize}

\end{frame}

\begin{frame}[fragile]{Forecasting in R}
\small

\begin{block}{R}
\begin{verbatim}
library(forecast)

model <- arima(x, order=c(1,0,1))

fc <- forecast(model, h=10)

plot(fc)
\end{verbatim}
\end{block}

\vspace{0.3cm}

\textbf{Output:}
\begin{itemize}
    \item Forecast values
    \item Confidence intervals
\end{itemize}

\end{frame}

\begin{frame}[fragile]{Forecasting in Python}
\small

\begin{block}{Python}
\begin{verbatim}
from statsmodels.tsa.arima.model import ARIMA

model = ARIMA(x, order=(1,0,1)).fit()

forecast = model.get_forecast(steps=10)

mean = forecast.predicted_mean
conf_int = forecast.conf_int()
\end{verbatim}
\end{block}

\vspace{0.3cm}

\textbf{Output:}
\begin{itemize}
    \item Forecast values
    \item Confidence intervals
\end{itemize}

\end{frame}

\begin{frame}{Forecasting: Summary}

\textbf{Steps:}
\begin{itemize}
    \item Fit a model (AR, MA, ARMA)
    \item Check residuals
    \item Generate forecasts
\end{itemize}

\vspace{0.4cm}

\textbf{Key idea:}
\begin{itemize}
    \item Good forecasts require a well-specified model
\end{itemize}

\vspace{0.4cm}

\textbf{Important:}
\begin{itemize}
    \item Forecast uncertainty increases over time
\end{itemize}

\end{frame}
\end{document}